\subsection{What a Matrix Is}

A matrix is a rectangular arrangement of mathematical objects, usually numbers.
That description is correct, but incomplete. A matrix may organise data, record
coefficients of equations, describe relations between quantities, or encode a rule
that transforms one vector into another. The same array can therefore play
very different roles in different problems.

\begin{definitionbox}[Matrix and its size]
A matrix with $m$ rows and $n$ columns is called an $m\times n$ matrix. It has
the form
\[
A=
\begin{pmatrix}
a_{11} & a_{12} & \cdots & a_{1n}\\
a_{21} & a_{22} & \cdots & a_{2n}\\
\vdots & \vdots & \ddots & \vdots\\
a_{m1} & a_{m2} & \cdots & a_{mn}
\end{pmatrix}.
\]
The entry $a_{ij}$ lies in row $i$ and column $j$.
\end{definitionbox}

\begin{interpretationbox}[How to read the indices]
The first index identifies the row and the second identifies the column. Thus
$a_{ij}$ is found by moving first to row $i$ and then to column $j$. This order
will matter later when matrices are multiplied and when systems of equations are
encoded.
\end{interpretationbox}

\begin{examplebox}[Reading a matrix before calculating]
Consider
\[
M=
\begin{pmatrix}
4 & -2 & 0\\
7 & 1 & 5
\end{pmatrix}.
\]
The matrix has two rows and three columns, so it is a $2\times3$ matrix. Its
entries include
\[
m_{11}=4,
\qquad
m_{12}=-2,
\qquad
m_{23}=5.
\]
No operation has been performed. We have only identified the shape and read the
positions correctly.
\end{examplebox}

\begin{definitionbox}[Square matrix]
A matrix is square when it has the same number of rows and columns. An
$n\times n$ matrix is called a square matrix of order $n$.
\end{definitionbox}

\begin{interpretationbox}[Shape is part of the mathematical meaning]
A row matrix, a column matrix, a square matrix, and a rectangular matrix may
contain similar entries but need not represent the same kind of object. A column
matrix may represent a vector, a square matrix may represent a transformation
from a space to itself, and a rectangular matrix may relate collections of
different sizes.
\end{interpretationbox}

\begin{summarybox}[Reading a matrix]
Before calculating, determine the number of rows and columns, identify the role
of the entries, and decide what the rows and columns represent. A matrix is not
understood merely by seeing its numbers; it is understood by reading its
structure.
\end{summarybox}
